\pdfoutput=1
\documentclass[10pt]{amsart}
\usepackage{mathtools,amssymb}
\DeclareMathAlphabet{\mathmod}{U}{eur}{m}{n}
\usepackage{enumitem,needspace,setspace}
\usepackage{tikz-cd}
\setlist[enumerate]{font=\normalfont}
\usepackage[a4paper,left=2.8cm,right=2.8cm,top=3cm,bottom=3cm]{geometry}
\usepackage[hidelinks]{hyperref}
\raggedbottom
\AddToHook{cmd/subsection/before}{\Needspace{4\baselineskip}}
\AddToHook{cmd/subsubsection/before}{\Needspace{6\baselineskip}}
\setstretch{1.1}
\setlength{\emergencystretch}{3em}
\usepackage{etoolbox}
\patchcmd{\abstract}{\leftmargin3pc}{\leftmargin\dimexpr2.5pc+0.4cm\relax}{}{\PackageWarning{diaz-park}{Abstract margin patch failed}}
\newtheorem{theorem}{Theorem}[section]
\newtheorem{proposition}[theorem]{Proposition}
\newtheorem{lemma}[theorem]{Lemma}
\newtheorem{corollary}[theorem]{Corollary}
\theoremstyle{definition}
\newtheorem{definition}[theorem]{Definition}
\newtheorem{remark}[theorem]{Remark}
\DeclareMathOperator{\Aut}{Aut}
\DeclareMathOperator{\Out}{Out}
\DeclareMathOperator{\Hom}{Hom}
\DeclareMathOperator{\Mor}{Mor}
\numberwithin{equation}{section}
\hypersetup{pdftitle={A proof of the sharpness conjectures for all saturated fusion systems},
pdfauthor={Baoyu Zhang and Shengmin Zhang},
pdfkeywords={saturated fusion system, p-local finite group, Mackey functor, centric orbit category, higher limit, homology decomposition}}
\allowdisplaybreaks
\renewcommand{\le}{\leqslant}
\renewcommand{\leq}{\leqslant}
\renewcommand{\ge}{\geqslant}
\renewcommand{\geq}{\geqslant}
\hyphenation{Theorem Corollary Proposition Section Definition Asch-bacher}
\newcommand{\F}{\mathcal F}
\newcommand{\Hh}{\mathcal H}
\newcommand{\Op}{\mathcal O}
\newcommand{\Inn}{\operatorname{Inn}}
\newcommand{\Iso}{\operatorname{Iso}}
\newcommand{\limi}[2]{\mathop{\varprojlim\nolimits^{#1}}_{#2}}
\title{A proof of the sharpness conjectures for all saturated fusion systems}
\author{Baoyu Zhang}
\address{University of Birmingham}
\email{baoyuzhang.math@outlook.com}
\author{Shengmin Zhang}
\address{Peking University}
\email{zhangshengmin26@stu.pku.edu.cn}
\date{}
\subjclass{Primary 20D20; Secondary 20J06, 55R35, 55R40, 18G10}
\keywords{Saturated fusion system, $p$-local finite group, Mackey functor, centric orbit category, higher limit, homology decomposition}
\begin{document}
\begin{abstract}
We prove the D\'iaz--Park conjecture for sharpness of saturated fusion systems. More generally, for every saturated fusion system $\F$ over a finite $p$-group and every Mackey functor for $\F$ over the localisation $\mathbb{Z}_{(p)}$ of the integers at $p$, the higher limits of its contravariant part over the centric orbit category vanish in all positive degrees. In particular, the subgroup and normaliser decompositions of every $p$-local finite group $(S,\F,\mathcal L)$ over the $\F$-centric subgroups are sharp. This answers in the affirmative a question in the list of important open problems in [M.~Aschbacher, R.~Kessar and B.~Oliver, \emph{Fusion Systems in Algebra and Topology}, Cambridge Univ.\ Press (2011)], and proves the sharpness conjecture stated in [M.~Aschbacher and B.~Oliver, \emph{Fusion systems}, Bull.\ Amer.\ Math.\ Soc.\ (2016)] and Conjecture~1.2 of [E.~Yal\c{c}\i{}n, \emph{Higher limits over the fusion orbit category}, Adv.\ Math.\ (2022)]. The proof is algebraic and does not rely on the classification of finite simple groups.
\end{abstract}
\maketitle
\section{Introduction}
Let $p$ be a prime. In 1976, Puig~\cite{Puig1976} studied the $p$-local structure of a finite group $G$ through the category whose morphisms are the homomorphisms between $p$-subgroups induced by conjugation in $G$. In a Bourbaki lecture~\cite{Puig1982}, he called the same category, with the nontrivial $p$-subgroups as objects, the \emph{Frobenius category} of $G$. He gave an abstract definition of Frobenius categories over finite $p$-groups in 1991~\cite[\S1.3]{Puig2006} and published this work later~\cite{Puig2006,Puig2009}. These categories are now called \emph{saturated fusion systems}. Broto, Levi and Oliver~\cite{BLO} gave an equivalent form of Puig's definition and introduced \emph{$p$-local finite groups} $(S,\F,\mathcal L)$, where $\mathcal L$ is a centric linking system associated to a saturated fusion system $\F$ over a finite $p$-group $S$. Let $\Op(\F^c)$ be the centric orbit category of $\F$, whose objects are the $\F$-centric subgroups of $S$. The geometric realisation $|\mathcal L|$ of $\mathcal L$ is homotopy equivalent to the homotopy colimit over $\Op(\F^c)$ of a functor that sends each $\F$-centric subgroup $P$ of $S$ to a space with the homotopy type of the classifying space $BP$ (see~\cite[Proposition~2.2]{BLO}). This is the \emph{subgroup decomposition} of $|\mathcal L|$. Libman~\cite{Libman2006} showed that $|\mathcal L|$ is also homotopy equivalent to the homotopy colimit of a functor on the poset of $\F$-conjugacy classes of nonempty chains of $\F$-centric subgroups, ordered by reverse inclusion of subchains. The functor sends the class of each chain to a space with the homotopy type of the classifying space of its automorphism group in $\mathcal L$. This is the \emph{normaliser decomposition} of $|\mathcal L|$.

For $n\ge0$, let $H^n(-;\mathbb{F}_p)$ be the \emph{cohomology functor} $P\mapsto H^n(P;\mathbb{F}_p)$ on $\Op(\F^c)$. The Bousfield\nobreakdash--Kan spectral sequence of the subgroup decomposition has $E_2^{i,n}=\limi{i}{\Op(\F^c)}H^n(-;\mathbb{F}_p)$ for $i,n\ge0$ and converges to $H^{i+n}(|\mathcal L|;\mathbb{F}_p)$. A \emph{homology decomposition} glues together classifying spaces of certain subgroups of a finite group $G$ into a homotopy colimit whose map to $BG$ induces isomorphisms in mod~$p$ homology. Dwyer~\cite[\S1]{Dwyer1998Sharp} called a homology decomposition \emph{sharp} when its spectral sequence collapses onto the vertical axis. For the subgroup decomposition of $|\mathcal L|$, this means that $E_2^{i,n}=0$ for all $i>0$ and $n\ge0$. By the stable elements theorem of Broto, Levi and Oliver~\cite[Theorem~5.8]{BLO}, $H^n(|\mathcal L|;\mathbb{F}_p)$ is the inverse limit of $H^n(-;\mathbb{F}_p)$ over $\Op(\F^c)$. The question whether every $p$-local finite group has a sharp subgroup decomposition appears in the list of important open problems in the book of Aschbacher, Kessar and Oliver~\cite[\S III.7]{AKO}, which grew out of the 2007 Birmingham workshop on fusion systems. As noted in the book, sharpness would give a new proof of the stable elements theorem. In the survey of Aschbacher and Oliver~\cite{AO}, the same vanishing appears as the \emph{sharpness conjecture} for $H^*(|\mathcal L|;\mathbb{F}_p)$.

Higher limits over $\Op(\F^c)$ also arise in the first problem of the same list~\cite[\S III.7, Problem~1]{AKO}, which concerns the existence and uniqueness of centric linking systems. Broto, Levi and Oliver~\cite[Proposition~3.1]{BLO} showed that the obstructions to existence and to uniqueness lie in higher limits of the functor $P\mapsto Z(P)$. Chermak~\cite{Chermak2013} solved this problem by direct constructions. Oliver~\cite{Oliver2013} adapted Chermak's argument to this obstruction theory and showed that these higher limits vanish in all degrees at least~$2$, and in all positive degrees when $p$ is odd. Glauberman and Lynd~\cite{GL2016} removed the dependence of Oliver's proof on the classification of finite simple groups.

For a finite group $G$, Dwyer~\cite[Theorem~10.3]{Dwyer1998Sharp} proved that for every $n\ge0$ the functor ${P\mapsto H^n(P;\mathbb{F}_p)}$ has vanishing higher limits in positive degrees on the centric orbit category of $G$, whose objects are the $p$-centric subgroups of $G$. Jackowski and McClure~\cite[Proposition~5.14]{JM1992} had shown earlier that the higher limits of every Mackey functor over $\mathbb{Z}_{(p)}$ vanish in positive degrees on the orbit category of all $p$-subgroups of $G$. D\'iaz and Park~\cite{DP} defined Mackey functors $M=(M^*,M_*)$ for $\F$, which we recall in Definition~\ref{def:mackey}. Here $M^*$ is a contravariant and $M_*$ a covariant functor on the orbit category $\Op(\F)$ of all subgroups of $S$, and the two are related by a Mackey formula. The higher limits of the restriction of $M^*$ to $\Op(\F^c)$ are denoted by $\limi{i}{\Op(\F^c)}M^*$ for $i\ge0$. For a saturated fusion system $\F$, D\'iaz and Park~\cite[Theorem~A]{DP} proved that $\limi{i}{\Op(\F^c)}M^*=0$ for all $i>0$ when $M$ is an $\F^c$-restricted Mackey functor over $\mathbb{Z}_{(p)}$, that is, a pair of functors on $\Op(\F^c)$ satisfying condition~(i) of Definition~\ref{def:mackey} and a Mackey formula in which the sum is taken only over double coset representatives whose corresponding intersections are $\F$-centric. Their proof combines a theorem of Jackowski and McClure~\cite[Corollary~5.16]{JM1992} with results of D\'iaz and Libman~\cite[Proposition~2.9 and Remark~2.10]{DL2009}. Motivated by Dwyer's theorem and by the vanishing theorem for $\F^c$-restricted Mackey functors, D\'iaz and Park~\cite[Conjecture]{DP} conjectured that $\limi{i}{\Op(\F^c)}M^*=0$ for every $i>0$, every saturated fusion system $\F$ over a finite $p$-group $S$ and every Mackey functor $M$ for $\F$ over $\mathbb{F}_p$ whose values are finite dimensional.

The conjecture does not follow from this vanishing theorem, since the restriction to $\Op(\F^c)$ of a Mackey functor for $\F$ need not be $\F^c$-restricted. By an example due to D\'iaz and Libman (see~\cite[Example~5.4]{DP}), this happens already for $H^1(-;\mathbb{F}_p)$ when $\F$ is realised by a finite $p$-group. Our main result is the following strengthening of the conjecture.

\begin{theorem}\label{thm:main}
Let $\F$ be a saturated fusion system over a finite $p$-group $S$ and let $M$ be a Mackey functor for $\F$ over $\mathbb{Z}_{(p)}$. Then $\limi{i}{\Op(\F^c)}M^*=0$ for every $i>0$.
\end{theorem}

Theorem~\ref{thm:main} also holds over every field $k$ of characteristic~$p$, since a Mackey functor over $k$ can be regarded as a Mackey functor over $\mathbb{Z}_{(p)}$ without changing its higher limits, as noted after~\eqref{eq:cochains}. For $k=\mathbb{F}_p$ and Mackey functors whose values are finite dimensional, this is the D\'iaz--Park conjecture. Each $H^n(-;\mathbb{F}_p)$ is the restriction to $\Op(\F^c)$ of the contravariant part of a Mackey functor for $\F$ over $\mathbb{F}_p$. Hence we obtain the following.

\begin{corollary}\label{cor:cohomology}
Let $\F$ be a saturated fusion system over a finite $p$-group $S$ and let $n\ge0$. Then $\limi{i}{\Op(\F^c)}H^n(-;\mathbb{F}_p)=0$ for every $i>0$.
\end{corollary}

By Corollary~\ref{cor:cohomology}, the question in the book of Aschbacher, Kessar and Oliver has a positive answer. The sharpness conjecture stated in the survey of Aschbacher and Oliver holds as well, and so does the conjecture stated by Yal\c{c}\i{}n~\cite[Conjecture~1.2]{Yalcin}. Glauberman and Lynd~\cite{GL} proved that the higher limits of $H^n(-;\mathbb{F}_p)$ vanish in positive degrees for every saturated fusion system when $n\le p-2$, using the simple Mackey composition factors of these functors. Yal\c{c}\i{}n~\cite[Theorem~1.8]{Yalcin} showed that it suffices to prove the sharpness of the subgroup decomposition for $p$-local finite groups whose fusion systems have a nontrivial centre. In the same paper, Yal\c{c}\i{}n~\cite[Theorem~1.9]{Yalcin} proved that the subgroup decomposition of a $p$-local finite group is sharp if and only if the normaliser decomposition of Libman~\cite{Libman2006} over the $\F$-centric subgroups is sharp. For a finite group $G$, the analogous equivalence for the subgroup and normaliser decompositions of $BG$ over the $p$-centric subgroups of $G$ follows from a theorem of Grodal and Smith~\cite[Theorem~1.1 and Remark~3.2]{GrodalSmith}. Grodal~\cite[Theorem~1.5]{Grodal} proved that these decompositions of $BG$ are sharp over the radical $p$-subgroups for the principal block, that is, the $p$-centric subgroups $P$ of $G$ with $O_p(N_G(P)/PC_G(P))=1$. His theorem concerns finite groups and this smaller collection, while our results concern every $p$-local finite group $(S,\F,\mathcal L)$ and the collection of all $\F$-centric subgroups. When $\F$ is the fusion system $\F_S(G)$ of a finite group $G$ with Sylow $p$-subgroup $S$, D\'iaz and Park~\cite[Theorem~B]{DP} proved their conjecture by translating the theorem of Jackowski and McClure~\cite[Corollary~5.16]{JM1992} to fusion systems, using a theorem of Grodal~\cite[Theorem~1.2]{Grodal}. In this case, Corollary~\ref{cor:cohomology} also follows from Dwyer's theorem. Indeed, by a lemma of Broto, Levi and Oliver~\cite[Lemma~1.3]{BLOInvent}, the higher limits of $H^n(-;\mathbb{F}_p)$ over $\Op(\F^c)$ are those over the centric orbit category of $G$.

A counterexample to the conjecture would therefore have been \emph{exotic}, that is, not of the form $\F_S(G)$ for any finite group $G$. D\'iaz and Park~\cite[Theorem~C]{DP} proved the conjecture when $S$ is nonabelian and has an abelian subgroup of index~$p$. These fusion systems include the exotic examples of Broto, Levi and Oliver~\cite[\S9]{BLO} and of Ruiz and Viruel~\cite{RV2004}. They also include the many exotic simple fusion systems classified by Oliver~\cite{Oliver2014}, together with Craven and Semeraro~\cite{COS2017} and then with Ruiz~\cite{OR2020}. Grazian and Marmo~\cite[Theorem~A]{GM} proved the conjecture for all saturated fusion systems over a Sylow $p$-subgroup of $\mathrm{G}_2(p)$ with $p\ge5$. These include the $27$ exotic fusion systems of Parker and Semeraro~\cite{PS2018} at $p=7$.

Henke, Libman and Lynd~\cite[Theorem~1.1]{HLL} proved that the subgroup decomposition of $|\mathcal L|$ is sharp when $\F$ has a punctured group. A \emph{punctured group}~\cite[Definition~2.16]{HLL} is a transporter system or locality whose objects are the nontrivial subgroups of $S$. By a theorem of Henke~\cite[Theorem~A]{Henke2019}, $\F$ has a punctured group when it is of characteristic $p$-type. Henke, Libman and Lynd~\cite[Theorem~5.2(c),(d)]{HLL} showed that all Parker--Stroth fusion systems~\cite{PS2015} are of characteristic $p$-type, as are the Clelland--Parker exotic fusion systems~\cite{CP2010} whose essential subgroups are abelian. Levi and Oliver~\cite{LO2002} constructed the Benson--Solomon fusion systems, which are exotic fusion systems at the prime~$2$, and Aschbacher and Chermak~\cite{AC2010} gave another construction, as fusion systems of amalgamated free products of finite groups. Solomon conjectured that these fusion systems are the only exotic simple fusion systems over $2$-groups (see~\cite[Conjecture~57.12]{Guido}). This conjecture also appears in the same list of important open problems~\cite[\S III.7]{AKO}, and Henke and Lynd~\cite{HL2022} have provided further evidence for it. However, as shown by Henke, Libman and Lynd~\cite[Theorem~1.4]{HLL}, $\F_{\mathrm{Sol}}(3)$ is, up to isomorphism, the only Benson--Solomon system with a punctured group.

Praderio Bova~\cite{BovaGreen} defined the Mackey algebra of a fusion system and proved a Green correspondence for $\F^c$-restricted Mackey functors. In the preprint~\cite[Theorem~C]{BovaAmalgams}, Praderio Bova stated the vanishing for the cohomology functors of all Clelland--Parker systems~\cite{CP2010} and all Parker--Stroth systems~\cite{PS2015}, and in another preprint~\cite[Theorem~D]{BovaInductive} asserted its extension to the polynomial fusion systems of Grazian, Parker, Semeraro and van Beek~\cite{GPSvB} and to all saturated fusion systems over $p$-groups of maximal class~\cite{GP2025} or over $p$-groups of $p$-rank~$2$ with $p$~odd~\cite{DRV2007,PS2019}. This assertion also includes the $16$ exotic fusion systems of van Beek~\cite{vanBeek2025} and the two Henke--Shpectorov exotic fusion systems~\cite{HenkeShpectorov}. Praderio Bova further stated that for every Mackey functor over $\mathbb{F}_p$ whose values are finite dimensional, the higher limits of degree at least~$2$ vanish for the Clelland--Parker and \mbox{Parker--Stroth} systems~\cite[\S4]{BovaAmalgams} and for the polynomial and Henke--Shpectorov fusion systems and for six of the van Beek fusion systems~\cite[\S4]{BovaInductive}. The case of degree~$1$ was left open in both preprints. In April 2026, Praderio Bova~\cite{BovaInductive} stated that the case of the cohomology functors was still open in general, for example for the Oliver--Ruiz exotic fusion systems~\cite[Theorem~A(c)]{OR2021}.

The vanishing of higher limits in the open cases mentioned in this introduction follows from Theorem~\ref{thm:main} and Corollary~\ref{cor:cohomology}, and our argument gives a uniform proof of the established cases of the sharpness conjectures for saturated fusion systems.
\section{Methods, background and preliminary results}
Throughout, $p$ is a prime and $\F$ is a fusion system over a finite $p$-group $S$. For fusion systems we follow the terminology of Aschbacher, Kessar and Oliver~\cite{AKO}.

\subsection{Organisation and methods of the proof}
The proof combines standard tools from finite group theory, algebraic topology and representation theory: centralisers and central quotients of saturated fusion systems, higher limits over orbit categories, and simple Mackey functors. It uses no case analysis of exotic fusion systems, no classifying space for $\F$ and no infinite group that realises $\F$, and it does not rely on the classification of finite simple groups.

We prove Theorem~\ref{thm:main} first over $\mathbb{F}_p$, by induction on $|S|$ over all saturated fusion systems. By a reduction of D\'iaz and Park~\cite[Proposition~4.3]{DP}, which we extend in Lemma~\ref{lem:reduction} to Mackey functors whose values need not be finite dimensional, it suffices to prove that $\limi{i}{\Op(\F^c)}M^*=0$ for every $i>0$ and every simple Mackey functor $M$ whose minimal subgroups (see Definition~\ref{def:simple-minimal}) are not $\F$-centric. Let $E$ be a nontrivial elementary abelian subgroup of $S$ that is fully centralised in $\F$. Figure~\ref{fig:induction} shows the induction step. The induction hypothesis must apply to all Mackey functors, not only the cohomology functors, since for a central subgroup $Z$ of $\F$ the functor $H^n(-;\mathbb{F}_p)$ need not factor through the functor $P\mapsto P/Z$ from $\Op(\F^c)$ to $\Op(\F/Z)$. To prove Proposition~\ref{prop:criterion}, we also need the vanishing of the inverse limit over $\Op(C_\F(E)^c)$ for each such $E$. This vanishing fails for $H^0(-;\mathbb{F}_p)$, but by Lemma~\ref{lem:degree-zero} it holds for the simple Mackey functors whose minimal subgroups are not $\F$-centric. The theorem over $\mathbb{Z}_{(p)}$ then follows by Nakayama's lemma and the finite generation of higher limits, first for Mackey functors whose values are finitely generated and then in general by Lemma~\ref{lem:closure}.

\begin{figure}[ht]
\centering
\begin{tikzcd}[column sep=2.2cm, row sep=2.2em, every label/.append style={font=\small, inner sep=0.8ex}]
\text{induction hypothesis}
  \arrow[d, Rightarrow, "{|C_S(E)/E|<|S|}"] \\
\begin{tabular}{@{}c@{}}Theorem~\ref{thm:main}\\over $\mathbb{F}_p$\\for $C_\F(E)/E$\end{tabular}
  \arrow[r, Rightarrow, "\text{Proposition~\ref{prop:central}}"]
& \begin{tabular}{@{}c@{}}Theorem~\ref{thm:main}\\over $\mathbb{F}_p$\\for $C_\F(E)$\end{tabular}
  \arrow[r, Rightarrow, "\text{Proposition~\ref{prop:criterion}}",
    "\text{for every $E$}"']
&[0.6cm] \begin{tabular}{@{}c@{}}Theorem~\ref{thm:main}\\over $\mathbb{F}_p$\\for $\F$\end{tabular}
\end{tikzcd}
\caption{The proof of Proposition~\ref{prop:criterion} uses Lemmas~\ref{lem:reduction}, \ref{lem:detection} and~\ref{lem:degree-zero}.}\label{fig:induction}
\end{figure}

Broto, Levi and Oliver~\cite[proof of Theorem~5.8]{BLO} proved the stable elements theorem by a related induction over central quotients and centralisers of elementary abelian subgroups. In that induction they proved the sharpness of their centraliser decomposition, on which the reduction of Yal\c{c}\i{}n~\cite[Theorem~1.8]{Yalcin} rests. For a $p$-local finite group $(S,\F,\mathcal L)$ with $S\ne1$, this decomposition~\cite[Theorem~2.6]{BLO} is a homotopy equivalence between $|\mathcal L|$ and the homotopy colimit of a contravariant functor on the full subcategory of $\F$ whose objects are the nontrivial elementary abelian subgroups of $S$ that are fully centralised in $\F$. The functor sends each such subgroup $E$ to a space with the homotopy type of the geometric realisation of the centric linking system that Broto, Levi and Oliver~\cite[Definition~2.4 and Proposition~2.5(c)]{BLO} constructed from $\mathcal L$ for the centraliser $C_\F(E)$. Their argument for this sharpness uses a cohomology decomposition theorem of Dwyer and Wilkerson~\cite[Theorem~1.2]{DW1992} and Lannes's theory of the $T$-functor~\cite{Lannes1992}. This functor is defined algebraically on unstable modules over the mod~$p$ Steenrod algebra. Lannes~\cite[Theorem~0.6]{Lannes1992} showed that, under suitable hypotheses, it computes the mod~$p$ cohomology of mapping spaces from classifying spaces of elementary abelian $p$-groups into $p$-completed spaces. Henke, Libman and Lynd~\cite[\S1.1]{HLL} wrote that the proof of the stable elements theorem by Broto, Levi and Oliver ``is indirect and requires heavy machinery such as Lannes's $T$\nobreakdash-functor theory'', and noted that this theorem is a direct consequence of the sharpness of the subgroup decomposition.

Our induction is much more direct and more elementary: it uses Lemma~\ref{lem:degree-zero} in place of the sharpness of the centraliser decomposition. The proof of Lemma~\ref{lem:degree-zero} involves neither Lannes's theory nor the theorem of Dwyer and Wilkerson. By Corollary~\ref{cor:cohomology} and an observation of Yal\c{c}\i{}n~\cite[\S9, p.~55]{Yalcin}, the centraliser decomposition is sharp, so our argument also gives a more direct proof of the sharpness that Broto, Levi and Oliver~\cite[proof of Theorem~5.8]{BLO} established.

\subsection{Fusion systems and higher limits}
Maps are written on the left and composed from right to left. For a morphism $\varphi\colon P\to R$ of $\F$ and a subgroup $Q\le P$, let $\varphi|_Q\colon Q\to\varphi(Q)$ be the isomorphism induced by $\varphi$. The set of isomorphisms in $\Hom_\F(P,R)$ for $P,R\le S$ is denoted by $\Iso_\F(P,R)$, and $\Out_\F(P)=\Aut_\F(P)/\Inn(P)$. For $x\in S$, the map $c_x\colon S\to S$ sends $y$ to $xyx^{-1}$, and ${}^xP=c_x(P)$ for $P\le S$.

\needspace{4\baselineskip}
\begin{definition}
For $P\le S$, let $P^\F$ be the set of $\F$-conjugates of $P$.
\begin{itemize}
\item A subgroup $P\le S$ is \emph{$\F$-centric} if $C_S(Q)\le Q$ for every $Q\in P^\F$.
\item The set of $\F$-centric subgroups of $S$ is denoted by $\F^c$, and the set of those $P\in\F^c$ with $O_p(\Out_\F(P))=1$ by $\F^{cr}$.
\item The \emph{orbit category} $\Op(\F)$ of $\F$ is the category whose objects are the subgroups of $S$ and whose morphism sets are $\Mor_{\Op(\F)}(P,R)=\Inn(R)\backslash\Hom_\F(P,R)$ for $P,R\le S$. For a morphism $\varphi$ of $\F$, we also write $\varphi$ for the morphism of $\Op(\F)$ that it represents.
\item For a set $\mathcal V$ of subgroups of $S$ that is closed under $\F$-conjugacy, the full subcategory of $\Op(\F)$ with object set $\mathcal V$ is denoted by $\Op^{\mathcal V}(\F)$. In particular, the \emph{centric orbit category} is $\Op(\F^c)\coloneqq\Op^{\F^c}(\F)$.
\end{itemize}
\end{definition}

\begin{definition}
Let $Z$ and $E$ be subgroups of $S$.
\begin{itemize}
\item The subgroup $Z$ is \emph{central} in $\F$ if $Z$ is normal in $S$ and, for all $P,R\le S$, each ${\varphi\in\Hom_\F(P,R)}$ extends to a morphism $\widehat\varphi\in\Hom_\F(PZ,RZ)$ such that $\widehat\varphi|_Z=\operatorname{Id}_Z$.
\item The \emph{centraliser} $C_\F(E)$ of $E$ in $\F$ is the fusion system over $C_S(E)$ whose morphisms ${P\to R}$, for $P,R\le C_S(E)$, are the morphisms $\varphi\in\Hom_\F(P,R)$ that extend to a morphism ${\widehat\varphi\in\Hom_\F(PE,RE)}$ such that $\widehat\varphi|_E=\operatorname{Id}_E$.
\end{itemize}
\end{definition}

Let $E\le S$ be abelian. Then $E$ is central in $C_\F(E)$. Suppose that $\F$ is saturated and $E$ is fully centralised in $\F$. Then $C_\F(E)$ is saturated, and Broto, Levi and Oliver~\cite[Proposition~2.5(a)]{BLO} showed that a subgroup $P\le C_S(E)$ is $C_\F(E)$-centric if and only if $E\le P$ and $P$ is $\F$-centric.

Let $A$ be a commutative ring. We write $A\text{-}\mathmod{mod}$ for the category of all $A$-modules. Let $\mathcal C$ be a small category. The \emph{higher limits} $\limi{n}{\mathcal C}$ for $n\ge0$ are the right derived functors of the inverse limit on the category of functors $\mathcal C^{\mathrm{op}}\to A\text{-}\mathmod{mod}$. Let $F\colon\mathcal C^{\mathrm{op}}\to A\text{-}\mathmod{mod}$ be a functor. Then $\limi{0}{\mathcal C}F=\varprojlim_{\mathcal C}F$. Often a functor $\Phi\colon\mathcal D\to\mathcal C$ is clear from the context, for example the inclusion of a subcategory. In that case, $\limi{n}{\mathcal D}F$ denotes the higher limits of $F\circ\Phi^{\mathrm{op}}$. By the dual of~\cite[Appendix~II, Proposition~3.3]{GZ1967}, the groups $\limi{n}{\mathcal C}F$ are the cohomology groups of the cochain complex induced by the bar resolution for $\mathcal C$~\cite[\S III.5.1]{AKO}. Its terms are
\begin{equation}\label{eq:cochains}
C^n(\mathcal C;F)=\prod_{c_0\to c_1\to\cdots\to c_n}F(c_0),
\end{equation}
where the product is taken over all composable sequences of $n$ morphisms in $\mathcal C$. The groups $C^n(\mathcal C;F)$ and the coboundary maps of this complex depend only on the underlying abelian groups and homomorphisms of $F$. Therefore the higher limits of $F$ are invariant under a change of rings: they do not change when the values of $F$ are regarded as $B$-modules via a homomorphism $B\to A$ of commutative rings, for example as abelian groups.

Jackowski, McClure and Oliver~\cite[Proposition~5.1(ii)]{JMO1992II} stated the following change of categories for functors to abelian groups. By the naturality of their isomorphism and the invariance of higher limits under a change of rings, it also holds for functors to $A$-modules.

\begin{lemma}[{Jackowski and McClure~\cite[Lemma~3.1]{JM1992}}]\label{lem:change}
Let $A$ be a commutative ring and let ${\Phi\colon\mathcal D\to\mathcal C}$ be a functor of small categories that has a left adjoint. Then, for every functor ${F\colon\mathcal C^{\mathrm{op}}\to A\text{-}\mathmod{mod}}$ and every $n\ge0$, there is an isomorphism of $A$-modules ${\limi{n}{\mathcal C}F\cong\limi{n}{\mathcal D}(F\circ\Phi^{\mathrm{op}})}$, natural in $F$.
\end{lemma}

\subsection{Mackey functors}\label{subsec:mackey}
D\'iaz and Park defined Mackey functors for fusion systems as follows.

\begin{definition}[{\cite[Definition~2.1]{DP}}]\label{def:mackey}
Let $A$ be a commutative ring. A \emph{Mackey functor for $\F$ over $A$} is a pair $M=(M^*,M_*)$ of functors $M^*\colon\Op(\F)^{\mathrm{op}}\to A\text{-}\mathmod{mod}$ and $M_*\colon\Op(\F)\to A\text{-}\mathmod{mod}$ that agree on objects and satisfy the following conditions, where $M(P)$ is their common value at $P\le S$ and $\iota^P_Q$ is the inclusion of a subgroup $Q\le P$.
\begin{enumerate}[label=\textup{(\roman*)}]
\item For every isomorphism $\varphi$ of $\F$, $M^*(\varphi)=M_*(\varphi^{-1})$.
\item For all subgroups $Q,R\le P\le S$,
\[
M^*(\iota^P_Q)\,M_*(\iota^P_R)=\sum_xM_*(\iota^Q_{Q\cap{}^xR})\,M^*(\iota^{{}^xR}_{Q\cap{}^xR})\,M_*(c_x|_R),
\]
where $x$ runs over a set of representatives of the double cosets $QxR$ in $P$.
\end{enumerate}
\end{definition}

Let $M$ be a Mackey functor for $\F$ over a commutative ring $A$. We use the following notation for the maps of $M$.
\begin{itemize}
\item For $Q\le P\le S$, the \emph{restriction} and \emph{transfer} maps are $\mathrm{r}^P_Q=M^*(\iota^P_Q)$ and $\mathrm{t}^P_Q=M_*(\iota^P_Q)$.
\item For an isomorphism $\varphi$ of $\F$, the \emph{conjugation} map is $\operatorname{iso}(\varphi)=M_*(\varphi)$.
\end{itemize}
In this notation, condition~(ii) of Definition~\ref{def:mackey} is the \emph{Mackey formula}
\begin{equation}\label{eq:mackey}
\mathrm{r}^P_Q\,\mathrm{t}^P_R=\sum_x\mathrm{t}^Q_{Q\cap{}^xR}\,\mathrm{r}^{{}^xR}_{Q\cap{}^xR}\operatorname{iso}(c_x|_R).
\end{equation}

These maps satisfy the following relations.
\begin{itemize}
\item (Identity) $\mathrm{r}^P_P=\mathrm{t}^P_P=\operatorname{iso}(c_x|_P)=\operatorname{Id}_{M(P)}$ for $x\in P\le S$.
\item (Transitivity) $\mathrm{r}^Q_R\,\mathrm{r}^P_Q=\mathrm{r}^P_R$ and $\mathrm{t}^P_Q\,\mathrm{t}^Q_R=\mathrm{t}^P_R$ for $R\le Q\le P\le S$, and $\operatorname{iso}(\psi)\operatorname{iso}(\varphi)=\operatorname{iso}(\psi\varphi)$ for isomorphisms $\varphi$ and $\psi$ of $\F$ such that $\psi\varphi$ is defined.
\item (Conjugation) $\operatorname{iso}(\varphi|_Q)\,\mathrm{r}^P_Q=\mathrm{r}^{\varphi(P)}_{\varphi(Q)}\operatorname{iso}(\varphi)$ and $\operatorname{iso}(\varphi)\,\mathrm{t}^P_Q=\mathrm{t}^{\varphi(P)}_{\varphi(Q)}\operatorname{iso}(\varphi|_Q)$ for every isomorphism $\varphi\colon P\to\varphi(P)$ of $\F$ and every $Q\le P$.
\end{itemize}
Thus for each $P\le S$ the group $\Aut_\F(P)$ acts on $M(P)$ via $\alpha\mapsto\operatorname{iso}(\alpha)$, and $M(P)$ is an $A\Out_\F(P)$\nobreakdash-module.

For $n\ge0$, the cohomology groups $H^n(P;\mathbb{F}_p)$ with trivial coefficients form a Mackey functor for $\F$ over $\mathbb{F}_p$ as $P$ ranges over the subgroups of $S$. Its restriction and transfer maps are those of group cohomology. For an isomorphism $\varphi$ of $\F$, the map $\operatorname{iso}(\varphi)$ is induced by $\varphi^{-1}$. The restriction of its contravariant part to $\Op(\F^c)$ is the cohomology functor $H^n(-;\mathbb{F}_p)$.

Let $\Hh$ be a fusion subsystem of $\F$ over a subgroup $T\le S$. The inclusion of $\Hh$ in $\F$ induces a functor $\Op(\Hh)\to\Op(\F)$. Restricting both parts of a Mackey functor $M$ for $\F$ along this functor gives a Mackey functor for $\Hh$ with values $M(P)$ for $P\le T$, since the maps $c_x|_R$ in \eqref{eq:mackey} with $x\in T$ are morphisms of $\Hh$. If every $\Hh$-centric subgroup is $\F$-centric, then the functor maps $\Op(\Hh^c)$ to $\Op(\F^c)$.

\begin{definition}
Let $M$ be a Mackey functor for $\F$ over a commutative ring $A$.
\begin{itemize}
\item A \emph{Mackey subfunctor} $N$ of $M$ is a collection of submodules $N(P)\le M(P)$, one for each ${P\le S}$, such that $\mathrm{r}^P_Q(N(P))\subseteq N(Q)$, $\mathrm{t}^P_Q(N(Q))\subseteq N(P)$ and $\operatorname{iso}(\varphi)(N(P))\subseteq N(\varphi(P))$ for all ${Q\le P\le S}$ and every isomorphism $\varphi\colon P\to\varphi(P)$ of $\F$.
\item A \emph{morphism} $M\to N$ of Mackey functors for $\F$ over $A$ is a collection of $A$-module homomorphisms $M(P)\to N(P)$, one for each ${P\le S}$, that commute with the restriction, transfer and conjugation maps.
\end{itemize}
\end{definition}

Every Mackey subfunctor $N$ of $M$, with the restricted maps, is a Mackey functor for $\F$ over $A$. The Mackey functors for $\F$ over $A$ form an abelian category. Its kernels, images and cokernels are constructed objectwise. Thus a sequence of Mackey functors is exact if and only if it is exact at every subgroup of $S$. By Lemma~\ref{lem:closure}(ii), it is enough to prove the vanishing of higher limits over $\Op(\F^c)$ for Mackey functors whose values are finitely generated.

\begin{samepage}
\begin{lemma}\label{lem:closure}
Let $A$ be a commutative ring and let $M$ be a Mackey functor for $\F$ over $A$. For $P\le S$, $x\in M(P)$ and $R\le S$, let $N_x(R)$ be the $A$-submodule of $M(R)$ generated by the elements $\mathrm{t}^R_U\operatorname{iso}(\alpha)\,\mathrm{r}^P_V(x)$ with $V\le P$, $U\le R$ and $\alpha\in\Iso_\F(V,U)$. Then the following hold.
\begin{enumerate}[label=\textup{(\roman*)}]
\item The modules $N_x(R)$ form a Mackey subfunctor $N_x$ of $M$ whose values are finitely generated, and $x\in N_x(P)$.
\item Let $n\ge0$. If $\limi{n}{\Op(\F^c)}N^*=0$ for every Mackey subfunctor $N$ of $M$ whose values are finitely generated, then $\limi{n}{\Op(\F^c)}M^*=0$.
\end{enumerate}
\end{lemma}
\end{samepage}

\begin{proof}
(i) As $S$ is finite, there are only finitely many triples $(V,U,\alpha)$, and so $N_x(R)$ is finitely generated. Taking $(V,U,\alpha)=(P,P,\operatorname{Id}_P)$, we obtain $x\in N_x(P)$. Since the maps of $M$ are $A$-linear, it remains to show that the restriction, transfer and conjugation maps send each generator $\mathrm{t}^R_U\operatorname{iso}(\alpha)\,\mathrm{r}^P_V(x)$ to a sum of such generators. Let $R\le Q\le S$ and let $\beta\colon R\to\beta(R)$ be an isomorphism of $\F$. By the transitivity and conjugation relations, $\mathrm{t}^Q_R\,\mathrm{t}^R_U=\mathrm{t}^Q_U$ and $\operatorname{iso}(\beta)\,\mathrm{t}^R_U\operatorname{iso}(\alpha)=\mathrm{t}^{\beta(R)}_{\beta(U)}\operatorname{iso}((\beta|_U)\alpha)$. Thus the transfer and conjugation maps send generators to generators. Let $W\le R$. For $y\in R$, set $\gamma_y=(c_y|_U)\alpha\colon V\to{}^yU$ and $V_y=\gamma_y^{-1}(W\cap{}^yU)$. By \eqref{eq:mackey} and the conjugation and transitivity relations,
\[
\mathrm{r}^R_W\,\mathrm{t}^R_U\operatorname{iso}(\alpha)\,\mathrm{r}^P_V=\sum_y\mathrm{t}^W_{W\cap{}^yU}\,\mathrm{r}^{{}^yU}_{W\cap{}^yU}\operatorname{iso}(\gamma_y)\,\mathrm{r}^P_V=\sum_y\mathrm{t}^W_{W\cap{}^yU}\operatorname{iso}(\gamma_y|_{V_y})\,\mathrm{r}^P_{V_y},
\]
where $y$ runs over a set of representatives of the double cosets $WyU$ in $R$. Thus $N_x$ is a Mackey subfunctor of $M$.

(ii) By (i), $M$ is the direct limit of its Mackey subfunctors whose values are finitely generated. Since $\Op(\F^c)$ has finitely many objects and morphisms, the proof of \cite[Proposition~III.5.5(b)]{AKO} applies without the noetherian hypothesis used in part~(a) of that proposition and shows that higher limits commute with this direct limit. The assertion follows.
\end{proof}

\begin{definition}\label{def:simple-minimal}
Let $M$ be a Mackey functor for $\F$ over a commutative ring $A$.
\begin{itemize}
\item The Mackey functor $M$ is \emph{simple} if $M\ne0$ and the only Mackey subfunctors of $M$ are $0$ and~$M$.
\item A subgroup $Q\le S$ is a \emph{minimal subgroup} for $M$ if $M(Q)\ne0$ and $M(P)=0$ for every proper subgroup $P$ of $Q$.
\end{itemize}
\end{definition}

\begin{lemma}\label{lem:simple}
Let $M$ be a simple Mackey functor for $\F$ over $\mathbb{F}_p$ and let $Q$ be a minimal subgroup for $M$. Then the following hold.
\begin{enumerate}[label=\textup{(\roman*)}]
\item The minimal subgroups for $M$ are the elements of $Q^\F$, and $M(L)$ is a simple $\mathbb{F}_p\Out_\F(L)$-module for every $L\in Q^\F$.
\item For every $U\le S$, if $x\in M(U)$ satisfies $\mathrm{r}^U_L(x)=0$ for every $L\in Q^\F$ with $L\le U$, then $x=0$. In particular, if no $\F$-conjugate of $Q$ is contained in $U$, then $M(U)=0$.
\item Suppose that $\F$ is saturated. If $L\in Q^\F$ and $T\le S$ satisfy $L\le T$ and $C_T(L)\not\le L$, then~$\mathrm{r}^T_L=0$.
\end{enumerate}
\end{lemma}

\begin{proof}
(i) D\'iaz and Park~\cite[Proposition~3.1]{DP} stated this result and noted that Webb's proof~\cite[Proposition~2.1]{Webb1993} for finite groups also applies to fusion systems.

(ii) By (i), $M(Q)$ is a simple $\mathbb{F}_p\Out_\F(Q)$-module, and $M$ is isomorphic by \cite[Proposition~3.2]{DP} to the simple Mackey functor that D\'iaz and Park~\cite[\S3]{DP} constructed from $Q$ and $M(Q)$. Let $U\le S$. By \cite[formula~(4)]{DP}, $M(U)$ is a direct sum taken over the $U$-conjugacy classes of the subgroups $L\in Q^\F$ contained in $U$, with a representative chosen in each class. For each chosen representative $L$, it follows from \cite[formula~(6)]{DP} that $\mathrm{r}^U_L$ is injective on the summand corresponding to $L$ and zero on the other summands. This proves (ii).

(iii) Over $\mathbb{F}_p$, Grazian and Marmo~\cite[Lemma~2.9]{GM} and D\'iaz and Park~\cite[Corollary~3.4]{DP} gave analogues for fusion systems of a result of Webb~\cite[Proposition~2.8]{Webb1993}. We use the decomposition of $M(T)$ in the proof of (ii), with $L$ chosen as the representative of its $T$-conjugacy class. By the lemma of Grazian and Marmo, the summand corresponding to $L$ is $0$, since $C_T(L)\not\le L$. As $\mathrm{r}^T_L$ is zero on the other summands, it follows that $\mathrm{r}^T_L=0$.
\end{proof}

\begin{lemma}\label{lem:reduction}
Suppose that $\F$ is saturated and that $\limi{i}{\Op(\F^c)}N^*=0$ for every $i>0$ and every simple Mackey functor $N$ for $\F$ over $\mathbb{F}_p$ whose minimal subgroups are not $\F$-centric. Then $\limi{i}{\Op(\F^c)}M^*=0$ for every $i>0$ and every Mackey functor $M$ for $\F$ over $\mathbb{F}_p$.
\end{lemma}

\begin{proof}
By Lemma~\ref{lem:closure}(ii), we may assume that the values of $M$ are finite dimensional. D\'iaz and Park~\cite[Proposition~4.3]{DP} showed that it suffices to prove that $\limi{i}{\Op(\F^c)}N^*=0$ for every $i>0$ and every simple composition factor $N$ of $M$ with a minimal subgroup that is not $\F$-centric. By Lemma~\ref{lem:simple}(i), all minimal subgroups for each such factor are $\F$-conjugate, so none is $\F$-centric and the assumed vanishing holds for these factors.
\end{proof}
\section{Quotients by central subgroups}
In this section $\F$ is saturated and $Z$ is a subgroup of $S$ that is central in $\F$. Then $Z\le Z(S)$, and every morphism $\varphi\colon P\to R$ of $\F$ between subgroups containing $Z$ is the identity on $Z$ and induces an injective homomorphism $\overline\varphi\colon P/Z\to R/Z$. The \emph{quotient} $\F/Z$ of $\F$ by $Z$~\cite[Definition~II.5.3]{AKO} is the fusion system over $S/Z$ whose morphisms are the homomorphisms $\overline\varphi$. By \cite[Lemma~II.5.5]{AKO}, the quotient $\F/Z$ is saturated. To compare higher limits over the centric orbit categories of $\F$ and $\F/Z$, we use the following set of subgroups of $S$.

\begin{lemma}[\cite{BCGLO}]\label{lem:quotient-collection}
Let $\mathcal X$ be the set of subgroups $P\le S$ that contain $Z$ and for which $P/Z$ is $\F/Z$-centric. Then the following hold.
\begin{enumerate}[label=\textup{(\roman*)}]
\item The set $\mathcal X$ is closed under $\F$-conjugacy and overgroups in $S$, and $\F^{cr}\subseteq\mathcal X\subseteq\F^c$.
\item The maps $P\mapsto P/Z$ on objects and $\varphi\mapsto\overline\varphi$ on representatives of morphisms define an isomorphism of categories from $\Op^{\mathcal X}(\F)$ to $\Op((\F/Z)^c)$.
\end{enumerate}
\end{lemma}

\begin{proof}
In (i) and (ii) we use arguments from the proof of the converse in~\cite[Corollary~6.14]{BCGLO}. That proof assumes that $\F/Z$ is the fusion system of a finite group, but these arguments require only that $\F$ is saturated and $Z$ is central in $\F$.

(i) As the set of $\F/Z$-centric subgroups of $S/Z$ is closed under $\F/Z$-conjugacy and overgroups, the set $\mathcal X$ is closed under $\F$-conjugacy and overgroups.
By a result of Broto, Castellana, Grodal, Levi and Oliver~\cite[Lemma~6.4(a)]{BCGLO}, we have $\mathcal X\subseteq\F^c$.
Their proof of~\cite[Corollary~6.14, p.~3856]{BCGLO} also shows that $O_p(\Out_\F(P))\ne1$ for every $\F$-centric subgroup $P$ outside $\mathcal X$.
Hence $\F^{cr}\subseteq\mathcal X$.

(ii) This isomorphism is established in~\cite[p.~3856, equation~(3)]{BCGLO}.
\end{proof}

\begin{proposition}\label{prop:central}
Suppose that $\limi{i}{\Op((\F/Z)^c)}N^*=0$ for every $i>0$ and every Mackey functor $N$ for $\F/Z$ over $\mathbb{F}_p$. Then $\limi{i}{\Op(\F^c)}M^*=0$ for every $i>0$ and every Mackey functor $M$ for $\F$ over~$\mathbb{F}_p$.
\end{proposition}

\begin{proof}
By Lemma~\ref{lem:reduction}, it suffices to show that $\limi{i}{\Op(\F^c)}M^*=0$ for every $i>0$ when $M$ is simple.
Let $M$ be simple and let $Q$ be a minimal subgroup for $M$.
As $Z$ is central in $\F$, either every element of $Q^\F$ contains $Z$, or none does.

Suppose first that $Z\not\le Q$, and let $P\in\F^c$.
Then $Z\le P$.
By Lemma~\ref{lem:simple}(iii), $\mathrm{r}^P_L=0$ for every $L\in Q^\F$ with $L\le P$, since $Z\le C_P(L)$ and $Z\not\le L$.
Thus $M(P)=0$ by Lemma~\ref{lem:simple}(ii).
Hence the restriction of $M^*$ to $\Op(\F^c)$ is zero, and so are its higher limits.

Suppose from now on that $Z\le Q$.
For $P\le S$ containing $Z$, let $\Gamma_P$ be the kernel of the natural homomorphism
\[
\Aut_\F(P)\longrightarrow\Aut_{\F/Z}(P/Z),
\]
which sends each $\F$-automorphism of $P$ to its induced automorphism of $P/Z$.
The kernel $\Gamma_P$ is normal in $\Aut_\F(P)$ and acts trivially on $Z$ and $P/Z$, so it is a $p$-group (see, for example, \cite[Corollary~5.3.3]{Gorenstein}).

Let $L\in Q^\F$.
Then $Z\le L$.
By Lemma~\ref{lem:simple}(i), the module $M(L)$ is a simple $\mathbb{F}_p\Out_\F(L)$-module.
As $\Gamma_L$ is a normal $p$-subgroup of $\Aut_\F(L)$, it acts trivially on $M(L)$.

Next, let $P\le S$ contain $Z$, let $\gamma\in\Gamma_P$ and let $L\in Q^\F$ with $L\le P$.
Since $L$ contains $Z$ and $\gamma$ induces the identity on $P/Z$, we have $\gamma(L)=L$ and $\gamma|_L\in\Gamma_L$.
As $\Gamma_L$ acts trivially on $M(L)$, the conjugation relations of Section~\ref{subsec:mackey} give $\mathrm{r}^P_L\operatorname{iso}(\gamma)=\operatorname{iso}(\gamma|_L)\,\mathrm{r}^P_L=\mathrm{r}^P_L$.
It follows from Lemma~\ref{lem:simple}(ii) that $\operatorname{iso}(\gamma)$ is the identity of $M(P)$.
As $\gamma^{-1}\in\Gamma_P$, both $M_*(\gamma)=\operatorname{iso}(\gamma)$ and ${M^*(\gamma)=M_*(\gamma^{-1})=\operatorname{iso}(\gamma^{-1})}$ are the identity of $M(P)$.

We now define a Mackey functor $N$ for $\F/Z$ over $\mathbb{F}_p$.
Set $N(P/Z)=M(P)$ for $Z\le P\le S$.
Every morphism of $\F/Z$ is $\overline\psi$ for some $\psi\in\Hom_\F(P,R)$ with $P$ and $R$ containing $Z$.
We set $N^*(\overline\psi)=M^*(\psi)$ and $N_*(\overline\psi)=M_*(\psi)$.
These maps do not depend on the choice of $\psi$.
Indeed, let $\chi\in\Hom_\F(P,R)$ also satisfy $\overline\chi=\overline\psi$.
Then $\psi(P)=\chi(P)$, as both contain $Z$ and have the same image in $R/Z$.
Thus $\chi=\psi\gamma$ for some $\gamma\in\Aut_\F(P)$.
As $\overline\psi\,\overline\gamma=\overline\psi$ and $\overline\psi$ is injective, we have $\gamma\in\Gamma_P$.
Since $M_*(\gamma)$ and $M^*(\gamma)$ are the identity of $M(P)$, it follows that $M_*(\chi)=M_*(\psi)M_*(\gamma)=M_*(\psi)$ and $M^*(\chi)=M^*(\gamma)M^*(\psi)=M^*(\psi)$.
Moreover, for $y\in R$ the morphism $c_{yZ}|_{R/Z}\,\overline\psi$ is induced by $c_y|_R\,\psi$, which represents the same morphism of $\Op(\F)$ as $\psi$.
Hence $N^*(\overline\psi)$ and $N_*(\overline\psi)$ depend only on the morphism of $\Op(\F/Z)$ represented by $\overline\psi$.

As $\psi\mapsto\overline\psi$ preserves composition and identities, $N^*$ and $N_*$ are functors on $\Op(\F/Z)^{\mathrm{op}}$ and on $\Op(\F/Z)$ that agree on objects.
If $\overline\psi$ is an isomorphism, then so is $\psi$, and $\overline{\psi^{-1}}=\overline\psi^{-1}$.
Condition~(i) of Definition~\ref{def:mackey} for $N$ therefore follows from condition~(i) for $M$.
For condition~(ii), let $U,W\le P\le S$ contain $Z$.
The restriction and transfer maps of $N$ are those of $M$.
For $x\in P$, the morphism $c_x|_W$ of $\F$ induces $c_{xZ}|_{W/Z}$, and so $N_*(c_{xZ}|_{W/Z})=M_*(c_x|_W)=\operatorname{iso}(c_x|_W)$.
Since $U$, $W$ and ${}^xW$ contain the central subgroup $Z$, the map $UxW\mapsto(U/Z)(xZ)(W/Z)$ is a bijection from the double cosets of $U$ and $W$ in $P$ onto those of $U/Z$ and $W/Z$ in $P/Z$.
For the same reason, $(U\cap{}^xW)/Z=(U/Z)\cap{}^{xZ}(W/Z)$.
Thus the Mackey formula~\eqref{eq:mackey} for $M$ and the subgroups $U,W\le P$ agrees term by term with~\eqref{eq:mackey} for $N$ and the subgroups $U/Z,W/Z\le P/Z$.
Therefore $N$ is a Mackey functor for $\F/Z$ over $\mathbb{F}_p$.

By Lemma~\ref{lem:quotient-collection}(i), the set $\mathcal X$ contains $\F^{cr}$ and is closed under $\F$-conjugacy and overgroups.
For such a set, Broto, Levi and Oliver~\cite[proof of Corollary~3.6]{BLO} showed, using only that $\F$ is saturated, that restriction from $\Op(\F^c)$ to $\Op^{\mathcal X}(\F)$ does not change the higher limits of functors $\Op(\F^c)^{\mathrm{op}}\to\mathbb{Z}_{(p)}\text{-}\mathmod{mod}$. By the invariance of higher limits under a change of rings, noted after~\eqref{eq:cochains}, restriction also preserves the higher limits of $M^*$.
Moreover, by the definition of $N$, the restriction of $M^*$ to $\Op^{\mathcal X}(\F)$ is the composite of $N^*$ with the isomorphism of categories of Lemma~\ref{lem:quotient-collection}(ii).
As $N$ is a Mackey functor for $\F/Z$ over $\mathbb{F}_p$, the hypothesis of the proposition applies to $N$ and we obtain
\[
\limi{i}{\Op(\F^c)}M^*\cong\limi{i}{\Op^{\mathcal X}(\F)}M^*\cong\limi{i}{\Op((\F/Z)^c)}N^*=0
\]
for every $i>0$.
\end{proof}
\section{Centralisers of elementary abelian subgroups}
In this section $\F$ is saturated and $\F^e$ is the full subcategory of $\F$ whose objects are the nontrivial elementary abelian subgroups of $S$ that are fully centralised in $\F$.

Let $A$ be a commutative ring and let $F\colon\Op(\F^c)^{\mathrm{op}}\to A\text{-}\mathmod{mod}$ be a functor. When $S\ne1$, Yal\c{c}\i{}n~\cite[Lemma~9.3 and Theorem~9.4]{Yalcin} constructed a spectral sequence
\[
E_2^{s,t}=\limi{s}{(\F^e)^{\mathrm{op}}}\bigl(E\mapsto\limi{t}{\Op(C_\F(E)^c)}F\bigr)\Longrightarrow\limi{s+t}{\Op(\F^c)}F,\qquad s,t\ge0.
\]
Under the hypothesis of Lemma~\ref{lem:detection} below, every term $E_2^{s,t}$ is zero, and so the lemma follows immediately from Yal\c{c}\i{}n's spectral sequence. That construction uses a discrete group that realises $\F$ and a space on which this group acts. We give an alternative algebraic proof that uses no such group. Instead, it uses pairs of subgroups as Broto, Levi and Oliver~\cite[proof of Theorem~2.6]{BLO} did, with morphisms of the orbit category in place of those of the linking system.

\begin{lemma}\label{lem:detection}
Let $A$ be a commutative ring and let $F\colon\Op(\F^c)^{\mathrm{op}}\to A\text{-}\mathmod{mod}$ be a functor. Suppose that $S\ne1$ and that $\limi{n}{\Op(C_\F(E)^c)}F=0$ for every $n\ge0$ and every $E$ in $\F^e$. Then $\limi{n}{\Op(\F^c)}F=0$ for every $n\ge0$.
\end{lemma}

\begin{proof}
In this proof, functors to $A\text{-}\mathmod{mod}$ on $\Op(\F^c)$ and on the categories $\mathcal P$, $\mathcal P_0$ and $\mathcal P_E$ below are contravariant, and functors on $\F^e$ and on the undercategories and overcategories below are covariant. We first replace $\Op(\F^c)$ and $F$ by a category $\mathcal P_0$ of pairs of subgroups and a functor $F_0$ with the same higher limits, and then show that the right Kan extension along a functor $\tau\colon\mathcal P_0^{\mathrm{op}}\to\F^e$ takes an injective resolution of $F_0$ to an injective resolution of the zero functor.

Let $\mathcal P$ be the category whose objects are the pairs $(P,E)$ with $P\in\F^c$ and $E$ a nontrivial elementary abelian subgroup of $Z(P)$, and whose morphisms $(P,E)\to(R,D)$ are the morphisms $\varphi\colon P\to R$ of $\Op(\F^c)$ with $D\le\varphi(E)$. This condition does not depend on the representative $\varphi$, because the representatives form an $\Inn(R)$-orbit and $D\le Z(R)$. Let $\pi\colon\mathcal P\to\Op(\F^c)$ be the functor $(P,E)\mapsto P$.

For $P\in\F^c$, the pair $\sigma(P)=(P,\Omega_1(Z(P)))$ is an object of $\mathcal P$ since $S\ne1$. Let $(R,D)$ be an object of $\mathcal P$ and let $\varphi\colon P\to R$ be a morphism of $\Op(\F^c)$. As $D$ centralises the $\F$-centric subgroup $\varphi(P)$, we have $D\le\varphi(\Omega_1(Z(P)))$. Thus $\Mor_{\mathcal P}(\sigma(P),(R,D))=\Mor_{\Op(\F^c)}(P,\pi(R,D))$, and $\sigma$ is a functor left adjoint to $\pi$. Applying Lemma~\ref{lem:change} to $\pi$, we obtain $\limi{n}{\Op(\F^c)}F\cong\limi{n}{\mathcal P}(F\circ\pi^{\mathrm{op}})$.

Let $\mathcal P_0$ be the full subcategory of $\mathcal P$ whose objects are the pairs $(P,E)$ with $E$ in $\F^e$, and let $j\colon\mathcal P_0\to\mathcal P$ be the inclusion. Given $(P,E)$ in $\mathcal P$, let $\rho\colon E\to E_0$ be an isomorphism of $\F$ such that $E_0$ is fully centralised in $\F$. By the extension property of saturated fusion systems, $\rho$ extends to a morphism $\widehat\rho\colon C_S(E)\to S$ of $\F$. Since $P\le C_S(E)$, the restriction $\widehat\rho|_P$ is an isomorphism from $(P,E)$ to the object $(\widehat\rho(P),E_0)$ of $\mathcal P_0$. Thus every object of $\mathcal P$ is isomorphic to an object of the full subcategory $\mathcal P_0$. Hence $j$ is an equivalence of categories, and so it has a left adjoint. Let $F_0=F\circ\pi^{\mathrm{op}}\circ j^{\mathrm{op}}$. By Lemma~\ref{lem:change}, we have $\limi{n}{\mathcal P}(F\circ\pi^{\mathrm{op}})\cong\limi{n}{\mathcal P_0}F_0$. It remains to show that $\limi{n}{\mathcal P_0}F_0=0$ for every $n\ge0$.

For a morphism $\varphi\colon(P,E)\to(R,D)$ of $\mathcal P_0$, let $\delta_\varphi\colon D\to E$ be the restriction of $\varphi^{-1}$ to $D$. As $D\le Z(R)$, the map $\delta_\varphi$ does not depend on the representative $\varphi$. Moreover, $\delta_{\psi\varphi}=\delta_\varphi\delta_\psi$ for every morphism $\psi$ of $\mathcal P_0$ with domain $(R,D)$. Thus $(P,E)\mapsto E$ and $\varphi\mapsto\delta_\varphi$ define a functor $\tau\colon\mathcal P_0^{\mathrm{op}}\to\F^e$. For $E$ in $\F^e$, let $\mathcal P_E$ be the subcategory of $\mathcal P_0$ whose objects are the pairs $(P,E)$ and whose morphisms are the $\varphi$ with $\varphi|_E=\operatorname{Id}_E$, and let $j_E\colon\mathcal P_E\to\mathcal P_0$ be the inclusion. By the definition of $C_\F(E)$ and the criterion of Broto, Levi and Oliver~\cite[Proposition~2.5(a)]{BLO} for its centric subgroups, the objects of $\mathcal P_E$ are the pairs $(P,E)$ with $P\in C_\F(E)^c$. The morphisms of $\mathcal P_E$ are the morphisms of $\Op(\F^c)$ represented by morphisms of $C_\F(E)$ between these subgroups. Therefore $(P,E)\mapsto P$ is an isomorphism from $\mathcal P_E$ to $\Op(C_\F(E)^c)$. Under this isomorphism, $F_0\circ j_E^{\mathrm{op}}$ corresponds to the restriction of $F$ to $\Op(C_\F(E)^c)$.

Let $E$ be in $\F^e$, let $a=(P,D)$ be an object of $\mathcal P_0$, let $\theta\colon E\to D$ be a morphism of $\F^e$ and set $U=\theta(E)$. As $E$ is fully centralised in $\F$, it follows from the extension property that the isomorphism $\theta^{-1}\colon U\to E$ extends to a morphism $\beta\colon C_S(U)\to S$ of $\F$. Since $U\le D\le Z(P)$, we have $P\le C_S(U)$ and $E=\beta(U)\le\beta(D)\le Z(\beta(P))$. Thus the restriction $\lambda=\beta|_P$ is a morphism from $a$ to the object $a_\theta=(\beta(P),E)$ of $\mathcal P_E$, and $\delta_\lambda=\theta$. Let $\lambda^{-1}$ be the inverse of $\lambda$ in $\Op(\F^c)$, and let $\psi\colon a\to b$ be a morphism of $\mathcal P_0$ with $b$ in $\mathcal P_E$ and $\delta_\psi=\theta$. Then $\psi|_U=\theta^{-1}=\lambda|_U$, and so the morphism $\chi=\psi\lambda^{-1}$ of $\Op(\F^c)$ is the identity on $E$. Thus $\chi$ is the unique morphism of $\mathcal P_E$ with $\psi=\chi\lambda$.

The functor $X\mapsto X\circ\tau$, from functors on $\F^e$ to functors on $\mathcal P_0$, has a right adjoint $\operatorname{Ran}_\tau$, given by right Kan extension along $\tau$~\cite[X.3, Theorem~1 and Corollary~2]{MacLane}. For a functor $Y$ on $\mathcal P_0$, the module $(\operatorname{Ran}_\tau Y)(E)$ is the inverse limit of the functor $(a,\theta)\mapsto Y(a)$ on the undercategory $E\downarrow\tau$. Its objects are the pairs $(a,\theta)$ with $a$ in $\mathcal P_0$ and $\theta\colon E\to\tau(a)$ a morphism of $\F^e$. The morphisms $(a,\theta)\to(c,\eta)$ are the morphisms $\psi\colon c\to a$ of $\mathcal P_0$ with ${\delta_\psi\theta=\eta}$. The functor $X\mapsto X\circ\tau$ is exact and preserves constant functors. Hence $\operatorname{Ran}_\tau$ sends injectives to injectives, and $\varprojlim_{(\F^e)^{\mathrm{op}}}\operatorname{Ran}_\tau Y\cong\varprojlim_{\mathcal P_0}Y$ naturally in $Y$.

By the factorisation $\psi=\chi\lambda$ above, the functor ${\mathcal P_E\to(E\downarrow\tau)^{\mathrm{op}}}$ given by $b\mapsto(b,\operatorname{Id}_E)$ has a left adjoint, which sends $(a,\theta)$ to $a_\theta$. Applying Lemma~\ref{lem:change} to this functor in degree~$0$, we obtain $(\operatorname{Ran}_\tau Y)(E)\cong\varprojlim_{\mathcal P_E}(Y\circ j_E^{\mathrm{op}})$ naturally in $Y$.

The functor $Y\mapsto Y\circ j_E^{\mathrm{op}}$ has a left adjoint $\operatorname{Lan}_{j_E^{\mathrm{op}}}$, given by left Kan extension along $j_E^{\mathrm{op}}$~\cite[X.3, formulas~(9) and~(10)]{MacLane}. For a functor $X$ on $\mathcal P_E$ and an object $a$ of $\mathcal P_0$, the module $(\operatorname{Lan}_{j_E^{\mathrm{op}}}X)(a)$ is the colimit of the functor $(b,\psi)\mapsto X(b)$ on the overcategory $j_E^{\mathrm{op}}\downarrow a$. Its objects are the pairs $(b,\psi)$ with $b$ in $\mathcal P_E$ and $\psi\colon a\to b$ a morphism of $\mathcal P_0$. The morphisms $(b,\psi)\to(c,\varphi)$ are the morphisms $\chi\colon c\to b$ of $\mathcal P_E$ with $\chi\varphi=\psi$. As $\delta_\chi=\operatorname{Id}_E$ for every morphism $\chi$ of $\mathcal P_E$, this category is the disjoint union of the full subcategories whose objects are the pairs $(b,\psi)$ with $\delta_\psi=\theta$, one for each morphism $\theta\colon E\to\tau(a)$ of $\F^e$. By the factorisation $\psi=\chi\lambda$ above, $(a_\theta,\lambda)$ is a final object of the subcategory for $\theta$. Hence $(\operatorname{Lan}_{j_E^{\mathrm{op}}}X)(a)\cong\bigoplus_\theta X(a_\theta)$, and the functor $\operatorname{Lan}_{j_E^{\mathrm{op}}}$ is exact. Since the functor $Y\mapsto Y\circ j_E^{\mathrm{op}}$ is right adjoint to this exact functor, it sends injectives to injectives.

Let $0\to F_0\to I^0\to I^1\to\cdots$ be an injective resolution of $F_0$ in the category of functors on $\mathcal P_0$, and let $I^\bullet$ be the complex $I^0\to I^1\to\cdots$. Set $J^n=\operatorname{Ran}_\tau I^n$ for $n\ge0$, and let $J^\bullet$ be the complex $J^0\to J^1\to\cdots$ of injective functors on $\F^e$. For $E$ in $\F^e$, the functor $Y\mapsto Y\circ j_E^{\mathrm{op}}$ is exact and sends injectives to injectives, and so $I^\bullet\circ j_E^{\mathrm{op}}$ is an injective resolution of $F_0\circ j_E^{\mathrm{op}}$. Since $J^\bullet(E)\cong\varprojlim_{\mathcal P_E}(I^\bullet\circ j_E^{\mathrm{op}})$, the cohomology of $J^\bullet(E)$ in degree~$n$ is isomorphic to $\limi{n}{\Op(C_\F(E)^c)}F$, and this group is zero for every $n\ge0$ by the hypothesis of the lemma. The sequence $0\to J^0\to J^1\to\cdots$ is therefore an injective resolution of the zero functor on $\F^e$. As the higher limits of the zero functor vanish, the complex $\varprojlim_{(\F^e)^{\mathrm{op}}}J^\bullet$ is exact. So is the isomorphic complex $\varprojlim_{\mathcal P_0}I^\bullet$, whose cohomology in degree~$n$ is $\limi{n}{\mathcal P_0}F_0$. This proves the lemma.
\end{proof}

The hypothesis of Lemma~\ref{lem:detection} includes degree~$0$, while the induction hypothesis in Section~\ref{sec:proof} concerns positive degrees only. In degree~$0$ we use the following lemma.

\begin{lemma}\label{lem:degree-zero}
Let $\Hh$ be a saturated fusion subsystem of $\F$ over a subgroup $T\le S$ such that every $\Hh$-centric subgroup is $\F$-centric, and let $M$ be a simple Mackey functor for $\F$ over $\mathbb{F}_p$ whose minimal subgroups are not $\F$-centric. Then $\varprojlim_{\Op(\Hh^c)}M^*=0$.
\end{lemma}

\begin{proof}
An element of $\varprojlim_{\Op(\Hh^c)}M^*$ is a family $(x_P)_{P\in\Hh^c}$ with $x_P\in M(P)$ and $M^*(\varphi)(x_R)=x_P$ for every morphism $\varphi\colon P\to R$ of $\Hh$ between $\Hh$-centric subgroups. Set $x=x_T$. Then $x_P=\mathrm{r}^T_P(x)$ for every $P\in\Hh^c$. It therefore suffices to show that $x=0$.

Let $P\in\Hh^c$, $L\le P$ and $\alpha\in\Aut_\Hh(P)$. As $\alpha\,\iota^P_L=\iota^P_{\alpha(L)}\,\alpha|_L$ and $M^*(\alpha)(x_P)=x_P$, we have
\[
M^*(\alpha|_L)\,\mathrm{r}^T_{\alpha(L)}(x)=M^*(\alpha|_L)\,\mathrm{r}^P_{\alpha(L)}(x_P)=\mathrm{r}^P_L\,M^*(\alpha)(x_P)=\mathrm{r}^T_L(x).
\]
This holds also when $L$ is not $\Hh$-centric, since $M^*$ is defined on all of $\Op(\F)$. Now let $L,K\le T$ and let $\beta\in\Iso_\Hh(L,K)$. By Alperin's fusion theorem~\cite[Theorem~4.52]{Craven} for the saturated fusion system $\Hh$, the isomorphism $\beta$ is a composite of restrictions of automorphisms of $\Hh$-centric subgroups of $T$. Applying the equation above to each of these restrictions in turn, we obtain $M^*(\beta)\,\mathrm{r}^T_K(x)=\mathrm{r}^T_L(x)$.

Let $Q$ be a minimal subgroup for $M$ and let $L\in Q^\F$ with $L\le T$. Since $L$ is not $\F$-centric and every $\Hh$-centric subgroup is $\F$-centric, $L$ is not $\Hh$-centric. Thus $C_T(L_0)\not\le L_0$ for some $L_0\in L^\Hh$. As ${L_0\in Q^\F}$, we have $\mathrm{r}^T_{L_0}=0$ by Lemma~\ref{lem:simple}(iii). For $\beta\in\Iso_\Hh(L,L_0)$, we obtain ${\mathrm{r}^T_L(x)=M^*(\beta)\,\mathrm{r}^T_{L_0}(x)=0}$. Therefore $x=0$ by Lemma~\ref{lem:simple}(ii) with $U=T$.
\end{proof}

\begin{proposition}\label{prop:criterion}
Suppose that $\limi{i}{\Op(C_\F(E)^c)}N^*=0$ for every $i>0$, every $E$ in $\F^e$ and every Mackey functor $N$ for $C_\F(E)$ over $\mathbb{F}_p$. Then $\limi{i}{\Op(\F^c)}M^*=0$ for every $i>0$ and every Mackey functor $M$ for $\F$ over $\mathbb{F}_p$.
\end{proposition}

\begin{proof}
By Lemma~\ref{lem:reduction}, it suffices to show that $\limi{i}{\Op(\F^c)}M^*=0$ for every $i>0$ when $M$ is a simple Mackey functor for $\F$ over $\mathbb{F}_p$ whose minimal subgroups are not $\F$-centric. If $S=1$, then there is no such $M$. Suppose that $S\ne1$. Let $M$ be such a functor, let $F$ be the restriction of $M^*$ to $\Op(\F^c)$ and let $E$ be in $\F^e$. The fusion system $C_\F(E)$ is saturated, and every $C_\F(E)$-centric subgroup is $\F$\nobreakdash-centric. By the choice of $M$, the hypotheses of Lemma~\ref{lem:degree-zero} therefore hold for $\Hh=C_\F(E)$, $T=C_S(E)$ and $M$, and so $\varprojlim_{\Op(C_\F(E)^c)}F=0$. The restriction $N$ of $M$ to $C_\F(E)$ is a Mackey functor for $C_\F(E)$ over $\mathbb{F}_p$, and $N^*$ and $F$ have the same restriction to $\Op(C_\F(E)^c)$. By the hypothesis of the proposition, for every $i>0$ we have $\limi{i}{\Op(C_\F(E)^c)}N^*=0$ and therefore $\limi{i}{\Op(C_\F(E)^c)}F=0$. Thus $F$ satisfies the hypotheses of Lemma~\ref{lem:detection} with $A=\mathbb{F}_p$, and by that lemma $\limi{i}{\Op(\F^c)}M^*=0$ for every $i\ge0$.
\end{proof}
\section{Proof of Theorem~\ref{thm:main}}\label{sec:proof}
\begin{proposition}\label{prop:field}
Let $\F$ be a saturated fusion system over a finite $p$-group $S$ and let $M$ be a Mackey functor for $\F$ over $\mathbb{F}_p$. Then $\limi{i}{\Op(\F^c)}M^*=0$ for every $i>0$.
\end{proposition}

\begin{proof}
We argue by induction on $|S|$ and assume that the proposition holds for every saturated fusion system over a $p$-group of order less than $|S|$.
By Proposition~\ref{prop:criterion}, it suffices to show that $\limi{i}{\Op(C_\F(E)^c)}N^*=0$ for every $i>0$, every $E$ in $\F^e$ and every Mackey functor $N$ for $C_\F(E)$ over $\mathbb{F}_p$.
Let $E$ be in $\F^e$ and set $\Hh=C_\F(E)$.
Then $\Hh$ is a saturated fusion system over $C_S(E)$ in which $E$ is central, and so the quotient $\Hh/E$ over $C_S(E)/E$ is saturated by~\cite[Lemma~II.5.5]{AKO}.
Since $E\ne1$, we have $|C_S(E)/E|<|S|$.
By the induction hypothesis for $\Hh/E$, we have $\limi{i}{\Op((\Hh/E)^c)}N^*=0$ for every $i>0$ and every Mackey functor $N$ for $\Hh/E$ over $\mathbb{F}_p$.
Applying Proposition~\ref{prop:central} with $C_S(E)$, $\Hh$ and $E$ in place of $S$, $\F$ and $Z$, we obtain $\limi{i}{\Op(\Hh^c)}N^*=0$ for every $i>0$ and every Mackey functor $N$ for $\Hh$ over $\mathbb{F}_p$.
\end{proof}

\begin{proof}[Proof of Theorem~\ref{thm:main}]
Let $M$ be a Mackey functor for $\F$ over $\mathbb{Z}_{(p)}$ and let $i>0$.
Suppose first that the values of $M$ are finitely generated.
Write $K$, $pM$ and $M/pM$ for the kernel, image and cokernel of multiplication by $p$ on $M$. We have exact sequences
\[
0\longrightarrow K\longrightarrow M\xrightarrow{p}pM\longrightarrow0
\qquad\text{and}\qquad
0\longrightarrow pM\longrightarrow M\longrightarrow M/pM\longrightarrow0.
\]
Since $p$ annihilates $K$ and $M/pM$, both are Mackey functors for $\F$ over $\mathbb{F}_p$.
By Proposition~\ref{prop:field}, the higher limits of $K^*$ and $(M/pM)^*$ over $\Op(\F^c)$ vanish in every positive degree.
They do not change when the values are regarded as $\mathbb{Z}_{(p)}$-modules, as noted after~\eqref{eq:cochains}.

Restricted to $\Op(\F^c)$, the contravariant parts of the two sequences are exact.
Their long exact sequences of higher limits contain the exact sequences
\[
\limi{i}{\Op(\F^c)}M^*\longrightarrow\limi{i}{\Op(\F^c)}(pM)^*\longrightarrow\limi{i+1}{\Op(\F^c)}K^*=0
\]
and
\[
\limi{i}{\Op(\F^c)}(pM)^*\longrightarrow\limi{i}{\Op(\F^c)}M^*\longrightarrow\limi{i}{\Op(\F^c)}(M/pM)^*=0.
\]
The first maps in these sequences are induced by multiplication by $p$ from $M$ to $pM$ and by the inclusion of $pM$ in $M$.
Both maps are surjective.
Their composite is induced by multiplication by $p$ on $M$, and so it is multiplication by $p$ on $\limi{i}{\Op(\F^c)}M^*$.
Thus $p\limi{i}{\Op(\F^c)}M^*=\limi{i}{\Op(\F^c)}M^*$.
As $\Op(\F^c)$ has finitely many objects and morphisms and the values of $M$ are finitely generated over the noetherian ring $\mathbb{Z}_{(p)}$, this module is finitely generated by \cite[Proposition~III.5.5(a)]{AKO}, and so it is zero by Nakayama's lemma.

In general, the first case applies to every Mackey subfunctor of $M$ whose values are finitely generated, and so $\limi{i}{\Op(\F^c)}M^*=0$ by Lemma~\ref{lem:closure}(ii) with $A=\mathbb{Z}_{(p)}$.
\end{proof}
\section*{Acknowledgements}
Shengmin Zhang first found a candidate proof of the D\'iaz--Park conjecture via interacting with ChatGPT. He then did an initial Lean formalisation conditional on published results, under Baoyu Zhang's direction. Shengmin Zhang invited Baoyu Zhang to write the paper together. Shengmin Zhang would like to thank his supervisors Prof.~Jiping Zhang and Prof.~Lizhong Wang for their support and encouragement in using AI.

Baoyu Zhang had formalised the book~\cite{Craven} and some fundamental papers of Aschbacher on fusion systems, and found which of the equivalent definitions of saturated fusion systems fits best with the other definitions and results on fusion systems in a formal development. To understand how far the candidate proof of the conjecture lay from the foundations of formal logic, Baoyu Zhang used Codex to formalise Theorem~\ref{thm:main} and Corollary~\ref{cor:cohomology} unconditionally in Lean and Rocq~\cite{ZZformal}. The formalisations contain the parts of the theory of fusion systems, category theory, representation theory and topology that they need and the libraries do not provide, and formal counterparts of the other numbered statements of this paper. They also prove topological consequences of Corollary~\ref{cor:cohomology}, such as the stable elements theorem. The Lean formal proofs of Theorem~\ref{thm:main} and Corollary~\ref{cor:cohomology} were checked with Nanoda, an independently implemented proof checker, and the Rocq formal proofs were rechecked with \texttt{coqchk}. Every result used in these formal proofs is formally proved, either in the formalisations or in the libraries, and only standard axioms are used. In addition, the checking code of \texttt{comparator}, a tool of the Lean developers, checked the Lean formal proof of Theorem~\ref{thm:main} against a copy of its formal statement in a separate file that imports the same formal definitions. Claude Code independently examined the logs and results of these checks. The project sources contain 61,849 lines of Lean code in 957 modules and 99,857 lines of Rocq code in 1,248 modules, not counting blank lines, comment lines or the libraries. With the compiled libraries reused, compiling the project sources took about 241 minutes of CPU time for Lean under Windows~11 and about 196 minutes for Rocq under Linux. The sampled peak memory of the compiler processes was about 7.1~GB for Lean and 10.5~GB for Rocq, and checking used at most about 2.2~GB with Nanoda and 7.9~GB with \texttt{coqchk}. Nanoda and \texttt{coqchk} ran under Linux and the checking code of \texttt{comparator} under Windows. The formalisation report in the same repository~\cite{ZZformal} describes the formalisations, compares them with the proofs in this paper and contains some corrections to the literature. It details these checks, lists Lean and Rocq declarations for each numbered statement, and explains in natural language the theory on which the proofs rest, the sources that the formalisations follow and how the results are traced down to the foundations of formal logic. Baoyu Zhang also completely rewrote the manuscript, simplified the proof significantly, wrote most of the explanations in the paper and greatly improved its exposition.
\begin{thebibliography}{99}

\bibitem{AC2010}
M. Aschbacher and A. Chermak,
{\it A group-theoretic approach to a family of $2$-local finite groups constructed by Levi and Oliver},
Ann.\ of Math.\ (2) \textbf{171} (2010), 881--978.

\bibitem{AKO}
M. Aschbacher, R. Kessar, and B. Oliver,
{\it Fusion Systems in Algebra and Topology},
London Math.\ Soc.\ Lecture Note Ser., vol.~391,
Cambridge Univ.\ Press, Cambridge, 2011.

\bibitem{AO}
M. Aschbacher and B. Oliver,
{\it Fusion systems},
Bull.\ Amer.\ Math.\ Soc.\ (N.S.) \textbf{53} (2016), 555--615.

\bibitem{BCGLO}
C. Broto, N. Castellana, J. Grodal, R. Levi, and B. Oliver,
{\it Extensions of $p$-local finite groups},
Trans.\ Amer.\ Math.\ Soc.\ \textbf{359} (2007), 3791--3858.

\bibitem{BLOInvent}
C. Broto, R. Levi, and B. Oliver,
{\it Homotopy equivalences of $p$-completed classifying spaces of finite groups},
Invent.\ Math.\ \textbf{151} (2003), 611--664.

\bibitem{BLO}
C. Broto, R. Levi, and B. Oliver,
{\it The homotopy theory of fusion systems},
J.\ Amer.\ Math.\ Soc.\ \textbf{16} (2003), 779--856.

\bibitem{Guido}
I. Chatterji (ed.),
{\it Guido's Book of Conjectures},
Monogr.\ Enseign.\ Math., vol.~40,
L'Enseignement Math\'ematique, Geneva, 2008.

\bibitem{Chermak2013}
A. Chermak,
{\it Fusion systems and localities},
Acta Math.\ \textbf{211} (2013), 47--139.

\bibitem{CP2010}
M. Clelland and C. Parker,
{\it Two families of exotic fusion systems},
J.\ Algebra \textbf{323} (2010), 287--304.

\bibitem{Craven}
D.~A. Craven,
{\it The Theory of Fusion Systems: An Algebraic Approach},
Cambridge Stud.\ Adv.\ Math., vol.~131,
Cambridge Univ.\ Press, Cambridge, 2011.

\bibitem{COS2017}
D.~A. Craven, B. Oliver, and J. Semeraro,
{\it Reduced fusion systems over $p$-groups with abelian subgroup of index~$p$:~II},
Adv.\ Math.\ \textbf{322} (2017), 201--268.

\bibitem{DL2009}
A. D\'iaz and A. Libman,
{\it The Burnside ring of fusion systems},
Adv.\ Math.\ \textbf{222} (2009), 1943--1963.

\bibitem{DP}
A. D\'iaz and S. Park,
{\it Mackey functors and sharpness for fusion systems},
Homology Homotopy Appl.\ \textbf{17} (2015), 147--164.

\bibitem{DRV2007}
A. D\'iaz, A. Ruiz, and A. Viruel,
{\it All $p$-local finite groups of rank two for odd prime $p$},
Trans.\ Amer.\ Math.\ Soc.\ \textbf{359} (2007), 1725--1764.

\bibitem{Dwyer1998Sharp}
W.~G. Dwyer,
{\it Sharp homology decompositions for classifying spaces of finite groups},
in {\it Group Representations: Cohomology, Group Actions and Topology},
Proc.\ Sympos.\ Pure Math., vol.~63, Amer.\ Math.\ Soc.,
Providence, RI, 1998, 197--220.

\bibitem{DW1992}
W.~G. Dwyer and C.~W. Wilkerson,
{\it A cohomology decomposition theorem},
Topology \textbf{31} (1992), 433--443.

\bibitem{GZ1967}
P. Gabriel and M. Zisman,
{\it Calculus of Fractions and Homotopy Theory},
Ergeb.\ Math.\ Grenzgeb., vol.~35, Springer, Berlin, 1967.

\bibitem{GL2016}
G. Glauberman and J. Lynd,
{\it Control of fixed points and existence and uniqueness of centric linking systems},
Invent.\ Math.\ \textbf{206} (2016), 441--484.

\bibitem{GL}
G. Glauberman and J. Lynd,
{\it Cohomology on the centric orbit category of a fusion system},
Fund.\ Math.\ \textbf{269} (2025), 187--199.

\bibitem{Gorenstein}
D. Gorenstein,
{\it Finite Groups}, 2nd ed.,
Chelsea, New York, 1980.

\bibitem{GM}
V. Grazian and E. Marmo,
{\it Sharpness of saturated fusion systems on a Sylow $p$-subgroup of
$\mathrm{G}_2(p)$},
Homology Homotopy Appl.\ \textbf{25} (2023), 329--342.

\bibitem{GP2025}
V. Grazian and C. Parker,
{\it Saturated fusion systems on $p$-groups of maximal class},
Mem.\ Amer.\ Math.\ Soc.\ \textbf{307} (2025), no.~1549.

\bibitem{GPSvB}
V. Grazian, C. Parker, J. Semeraro, and M. van Beek,
{\it Fusion systems related to polynomial representations of
$\mathrm{SL}_2(q)$},
J.\ London Math.\ Soc.\ (2) \textbf{113} (2026), Art.~e70481.

\bibitem{Grodal}
J. Grodal,
{\it Higher limits via subgroup complexes},
Ann.\ of Math.\ (2) \textbf{155} (2002), 405--457.

\bibitem{GrodalSmith}
J. Grodal and S.~D. Smith,
{\it Propagating sharp group homology decompositions},
Adv.\ Math.\ \textbf{200} (2006), 525--538.

\bibitem{Henke2019}
E. Henke,
{\it Subcentric linking systems},
Trans.\ Amer.\ Math.\ Soc.\ \textbf{371} (2019), 3325--3373.

\bibitem{HLL}
E. Henke, A. Libman, and J. Lynd,
{\it Punctured groups for exotic fusion systems},
Trans.\ London Math.\ Soc.\ \textbf{10} (2023), 21--99.

\bibitem{HL2022}
E. Henke and J. Lynd,
{\it Fusion systems with Benson--Solomon components},
Duke Math.\ J.\ \textbf{171} (2022), 673--737.

\bibitem{HenkeShpectorov}
E. Henke and S. Shpectorov,
{\it Fusion systems over Sylow $p$-subgroups of $\mathrm{PSp}_4(p^a)$},
unpublished manuscript.

\bibitem{JM1992}
S. Jackowski and J. McClure,
{\it Homotopy decomposition of classifying spaces via elementary abelian
subgroups},
Topology \textbf{31} (1992), 113--132.

\bibitem{JMO1992II}
S. Jackowski, J. McClure, and B. Oliver,
{\it Homotopy classification of self-maps of $BG$ via $G$-actions. Part II},
Ann.\ of Math.\ (2) \textbf{135} (1992), 227--270.

\bibitem{Lannes1992}
J. Lannes,
{\it Sur les espaces fonctionnels dont la source est le classifiant d'un $p$-groupe ab\'elien \'el\'ementaire},
Inst.\ Hautes \'Etudes Sci.\ Publ.\ Math.\ \textbf{75} (1992), 135--244.

\bibitem{LO2002}
R. Levi and B. Oliver,
{\it Construction of $2$-local finite groups of a type studied by Solomon and Benson},
Geom.\ Topol.\ \textbf{6} (2002), 917--990; and
{\it Correction to: Construction of $2$-local finite groups of a type studied by Solomon and Benson},
Geom.\ Topol.\ \textbf{9} (2005), 2395--2415.

\bibitem{Libman2006}
A. Libman,
{\it The normaliser decomposition for $p$-local finite groups},
Algebr.\ Geom.\ Topol.\ \textbf{6} (2006), 1267--1288.

\bibitem{MacLane}
S. Mac Lane,
{\it Categories for the Working Mathematician}, 2nd ed.,
Grad.\ Texts in Math., vol.~5,
Springer, New York, 1998.

\bibitem{Oliver2013}
B. Oliver,
{\it Existence and uniqueness of linking systems: Chermak's proof via obstruction theory},
Acta Math.\ \textbf{211} (2013), 141--175.

\bibitem{Oliver2014}
B. Oliver,
{\it Simple fusion systems over $p$-groups with abelian subgroup of index~$p$: I},
J.\ Algebra \textbf{398} (2014), 527--541.

\bibitem{OR2020}
B. Oliver and A. Ruiz,
{\it Reduced fusion systems over $p$-groups with abelian subgroup of index~$p$: III},
Proc.\ Roy.\ Soc.\ Edinburgh Sect.~A \textbf{150} (2020), 1187--1239.

\bibitem{OR2021}
B. Oliver and A. Ruiz,
{\it Simplicity of fusion systems of finite simple groups},
Trans.\ Amer.\ Math.\ Soc.\ \textbf{374} (2021), 7743--7777.

\bibitem{PS2018}
C. Parker and J. Semeraro,
{\it Fusion systems over a Sylow $p$-subgroup of $\mathrm{G}_2(p)$},
Math.\ Z.\ \textbf{289} (2018), 629--662.

\bibitem{PS2019}
C. Parker and J. Semeraro,
{\it Fusion systems on maximal class $3$-groups of rank two revisited},
Proc.\ Amer.\ Math.\ Soc.\ \textbf{147} (2019), 3773--3786.

\bibitem{PS2015}
C. Parker and G. Stroth,
{\it A family of fusion systems related to the groups $\mathrm{Sp}_4(p^a)$ and $\mathrm{G}_2(p^a)$},
Arch.\ Math.\ (Basel) \textbf{104} (2015), 311--323.

\bibitem{BovaGreen}
M. Praderio Bova,
{\it Green correspondence on centric Mackey functors over fusion systems},
J.\ Algebra \textbf{638} (2024), 279--357.

\bibitem{BovaAmalgams}
M. Praderio Bova,
{\it Computing higher limits over the fusion orbit category via amalgams},
arXiv:2601.14983v1, 21~January 2026.

\bibitem{BovaInductive}
M. Praderio Bova,
{\it An inductive approach to the D\'iaz--Park sharpness conjecture},
arXiv:2604.21161v1, 23~April 2026.

\bibitem{Puig1976}
L. Puig,
{\it Structure locale dans les groupes finis},
M\'em.\ Soc.\ Math.\ France \textbf{47} (1976), 5--132.

\bibitem{Puig1982}
L. Puig,
{\it La classification des groupes finis simples: bref aper\c{c}u et quelques cons\'equences internes},
S\'eminaire Bourbaki, exp.~no.~584,
Ast\'erisque \textbf{92--93} (1982), 101--128.

\bibitem{Puig2006}
L. Puig,
{\it Frobenius categories},
J.\ Algebra \textbf{303} (2006), 309--357.

\bibitem{Puig2009}
L. Puig,
{\it Frobenius Categories versus Brauer Blocks: The Grothendieck Group of the Frobenius Category of a Brauer Block},
Progr.\ Math., vol.~274,
Birkh\"auser, Basel, 2009.

\bibitem{RV2004}
A. Ruiz and A. Viruel,
{\it The classification of $p$-local finite groups over the extraspecial group of order $p^3$ and exponent $p$},
Math.\ Z.\ \textbf{248} (2004), 45--65.

\bibitem{vanBeek2025}
M. van Beek,
{\it Exotic fusion systems related to sporadic simple groups},
Math.\ Proc.\ Cambridge Philos.\ Soc.\ \textbf{178} (2025), 331--384.

\bibitem{Webb1993}
P. Webb,
{\it Two classifications of simple Mackey functors with applications to group cohomology and the decomposition of classifying spaces},
J.\ Pure Appl.\ Algebra \textbf{88} (1993), 265--304.

\bibitem{Yalcin}
E. Yal\c{c}\i{}n,
{\it Higher limits over the fusion orbit category},
Adv.\ Math.\ \textbf{406} (2022), Art.~108482.

\bibitem{ZZformal}
B. Zhang,
{\it Lean and Rocq formalisations of the sharpness conjectures for saturated fusion systems},\newline
\url{https://github.com/BaoyuZhang747/sharpness-Lean-Rocq}, 2026.

\end{thebibliography}
\end{document}
